\chapter{Εισαγωγή}

Η παρούσα αναφορά τεκμηριώνει τη μεθοδολογία, τις αρχιτεκτονικές επιλογές
και τα ευρήματα ενός συστήματος ανίχνευσης audiovisual deepfake, βασισμένου
σε late fusion παγωμένων χαρακτηριστικών ομιλίας και οπτικής κίνησης χειλιών,
πάνω σε ένα υποσύνολο του συνόλου δεδομένων MAVOS-DD.
Η εργασία ακολουθεί φασική δομή υλοποίησης (Phases 1–6)
και δίνει έμφαση στην ειλικρινή αξιολόγηση: τι μαθαίνει πράγματι
το μοντέλο, ποια shortcuts παρατηρούνται, και πού θα απαιτούνταν
αρχιτεκτονικές επεκτάσεις για αληθινή πολυτροπική κατανόηση.

% ─────────────────────────────────────────────────────────────────────────────
\section{Πρόβλημα και Κίνητρο}
% ─────────────────────────────────────────────────────────────────────────────

Το αντικείμενο της εργασίας είναι η \emph{πολυτροπική ανίχνευση deepfake} σε
βίντεο: διαχωρισμός αυθεντικού οπτικοακουστικού υλικού από συνθετικό,
ανεξαρτήτως ποιο από τα δύο κανάλια έχει παραποιηθεί. Το πρόβλημα
διαφέρει ουσιαστικά από την κλασική αντι-spoofing ομιλίας
(audio-only anti-spoofing), η οποία εστιάζει αποκλειστικά
στην αυθεντικότητα του ηχητικού σήματος.

Στο σύνολο δεδομένων MAVOS-DD που χρησιμοποιούμε, οι πηγές
παραποίησης εμπίπτουν σε δύο διακριτές κατηγορίες:

\begin{itemize}
  \item \textbf{Οπτικές παραποιήσεις (visual-fake):} Οι γεννήτριες
    \codeid{EchoMimic}, \codeid{MEMO}, \codeid{LivePortrait} και
    \codeid{Sonic} κινούν συνθετικά ένα πρόσωπο χρησιμοποιώντας
    \emph{πραγματικό ανθρώπινο ήχο} ως οδηγό.
    Το ηχητικό κανάλι του αρχείου \codeid{.mp4} παραμένει αυθεντικό·
    η παραποίηση εντοπίζεται αποκλειστικά στον οπτικό δρόμο.
  \item \textbf{Ηχητικές παραποιήσεις (audio-fake / TTS):} Οι μηχανές
    σύνθεσης ομιλίας \engineName{ElevenLabs}, \engineName{Google Neural2}
    και \engineName{OpenAI TTS} παράγουν ηχητικό αρχείο από μεταγραφή
    του αυθεντικού ομιλητή. Το βίντεο παραμένει αυθεντικό· η παραποίηση
    εντοπίζεται αποκλειστικά στο ηχητικό κανάλι.
\end{itemize}

Αυθεντικά βίντεο (\emph{bonafide}) περιέχουν πραγματικό πρόσωπο και
πραγματικό ήχο, και τα δύο από YouTube αρχεία πηγής.

Η δυαδική ετικέτα είναι: πραγματικό (real) → 0, ψεύτικο (fake) → 1,
ανεξαρτήτως κατηγορίας παραποίησης.
Ο στόχος είναι ένας ενιαίος ταξινομητής που να ανταποκρίνεται
σε οποιαδήποτε συνδυαστική επίθεση, αξιοποιώντας και τα δύο
κανάλια πληροφορίας.

% ─────────────────────────────────────────────────────────────────────────────
\section{Ερευνητικά Ερωτήματα}
% ─────────────────────────────────────────────────────────────────────────────

Η εργασία δεν θέτει ως στόχο την επίτευξη υψηλής ακρίβειας
με κάθε κόστος. Αντίθετα, ο πυρήνας της ανάλυσης οργανώνεται γύρω
από τέσσερα ψευδαισθητά ερωτήματα:

\begin{enumerate}

  \item \textbf{Τι μαθαίνει στην πραγματικότητα το μοντέλο σε αυτό το dataset;}
    Η υψηλή απόδοση σε επικύρωση δεν συνεπάγεται ότι το μοντέλο
    έχει αποκτήσει κατανόηση της αυθεντικότητας του ομιλητή.
    Εξετάζουμε εάν οι βελτιώσεις οφείλονται σε βαθύτερα ακουστικά
    ή οπτικά χαρακτηριστικά, ή σε επιφανειακά shortcuts
    (κωδικοποίηση αρχείου, φασματική ανομοιότητα γεννήτριας κ.ά.).

  \item \textbf{Είναι το audio anti-spoof ένας ανιχνευτής engine fingerprint;}
    Μετά τη διόρθωση του codec mismatch και της voice leakage,
    τα μοντέλα \codeid{WavLM} και \codeid{HuBERT} κορεσμένα
    στο \AUC{}~= 1.0 στο validation set υποδηλώνουν ότι
    η κυρίαρχη πληροφορία δεν είναι η αυθεντικότητα της ομιλίας,
    αλλά τα φασματικά αποτυπώματα κάθε TTS μηχανής.
    Αξιολογούμε αν ένα νέο, άγνωστο TTS σύστημα
    θα αποφύγει την ανίχνευση.

  \item \textbf{Μπορεί το visual branch να διακρίνει deepfake πρόσωπα από
    αυθεντικά όταν τα lip features δεν είναι πανομοιότυπα;}
    Στα Phases 1–5, κάθε bonafide γραμμή και η αντίστοιχη spoof γραμμή
    της μοιράζονταν τον ίδιο \codeid{source\_video\_id}
    και κατ' επέκταση τα ίδια \codeid{.npz} χαρακτηριστικά χειλιών.
    Ένα BiGRU που καλείται να διαχωρίσει δύο γραμμές με πανομοιότυπη
    είσοδο δεν μπορεί να ξεπεράσει τη συχνότητα κλάσης.
    Στo Phase 6+ οι πηγές \codeid{LivePortrait} και \codeid{Sonic}
    εισάγουν πρόσωπα όπου τα ίδια τα χαρακτηριστικά χειλιών διαφέρουν
    μεταξύ πραγματικού και ψεύτικου, επιτρέποντας ουσιαστική
    οπτική εκμάθηση.

  \item \textbf{Προσθέτει κάτι η συγχώνευση πέρα από το dominant modality;}
    Αν το οπτικό branch παραμένει τυχαίο, η concat-fusion
    κληρονομεί απλώς τη βαθμολογία του ηχητικού branch.
    Εξετάζουμε κατά πόσο, με λειτουργικά διαφορετικά οπτικά
    χαρακτηριστικά (Phase 6+), η late fusion ξεπερνά
    αμφότερες τις μεμονωμένες μονότροπες αρχιτεκτονικές.

\end{enumerate}

% ─────────────────────────────────────────────────────────────────────────────
\section{Εργαλεία και Τεχνολογίες}
% ─────────────────────────────────────────────────────────────────────────────

Η υλοποίηση βασίζεται αποκλειστικά σε ανοιχτού κώδικα ή προσβάσιμα
εργαλεία, με Python 3.10 ως κεντρική γλώσσα:

\begin{itemize}
  \item \textbf{Βαθιά Εκμάθηση:} \codeid{PyTorch}
    για τον ορισμό, εκπαίδευση και αξιολόγηση όλων των νευρωνικών
    μοντέλων· \codeid{scikit-learn} για βοηθητικές λειτουργίες
    (stratified splits, μετρικές αξιολόγησης) — δεν αποτελεί
    τον τελικό ταξινομητή.

  \item \textbf{Παγωμένα ηχητικά backbones (Hugging Face Transformers):}
    Τρία προ-εκπαιδευμένα μοντέλα, όλα 768-διάστατης εξόδου,
    χρησιμοποιούνται κατεψυγμένα — δεν εκπαιδεύονται ξανά:
    \begin{itemize}
      \item \codeid{Wav2Vec2} (\codeid{facebook/wav2vec2-base-960h})
      \item \codeid{WavLM} (\codeid{microsoft/wavlm-base})
      \item \codeid{HuBERT} (\codeid{facebook/hubert-base-ls960})
    \end{itemize}
    Κάθε backbone εξάγει χαρακτηριστικά μία φορά και τα αποθηκεύει
    σε \codeid{.npy} αρχεία· τα ακατέργαστα βίντεο δεν εισέρχονται
    ποτέ στον βρόχο εκπαίδευσης.

  \item \textbf{Οπτική εξαγωγή χαρακτηριστικών:}
    \codeid{MediaPipe FaceMesh} για την ανίχνευση
    40 σημείων αναφοράς χειλιών ανά καρέ. Αποτυχίες ανίχνευσης
    προσώπου καταγράφονται μέσω μάσκας στο \codeid{.npz} αρχείο·
    ποτέ δεν αποκρύπτονται αθόρυβα.

  \item \textbf{Επεξεργασία πολυμέσων:}
    \codeid{ffmpeg} με \codeid{libmp3lame} για αποκωδικοποίηση
    βίντεο, εξαγωγή ήχου, και κυκλική κωδικοποίηση MP3
    (codec mismatch neutralization).

  \item \textbf{Μεταγραφή ομιλίας:} Google STT V2 για την παραγωγή
    μεταγραφών από αυθεντικά αρχεία WAV, οι οποίες στη συνέχεια
    τροφοδοτούν τις μηχανές TTS.

  \item \textbf{Μηχανές σύνθεσης ομιλίας (TTS):}
    \engineName{ElevenLabs TTS/STS},
    \engineName{Google Neural2 TTS},
    \engineName{OpenAI TTS} (Phase 6+),
    \codeid{Coqui XTTS-v2} (τοπικά, Phase 6+ multi-engine).

  \item \textbf{Σύνολο δεδομένων:} \codeid{unibuc-cs/MAVOS-DD}
    από το Hugging Face Hub. Αγγλόφωνο υποσύνολο, με ανώτατα
    όρια (caps) ανά πηγή (Phases 1–5: 1.000 βίντεο·
    Phase 6+: περίπου 4.149 βίντεο).
\end{itemize}

% ─────────────────────────────────────────────────────────────────────────────
\section{Δομή της Αναφοράς}
% ─────────────────────────────────────────────────────────────────────────────

Η αναφορά ακολουθεί τη φασική πρόοδο της εργασίας:

\textbf{Κεφάλαιο~2} παρουσιάζει το σύνολο δεδομένων MAVOS-DD,
τις πηγές spoof και τη διαδικασία δειγματοληψίας υποσυνόλου.
\textbf{Κεφάλαιο~3} καλύπτει την προεπεξεργασία:
εξαγωγή ήχου, παραθύρωση, εξαγωγή χαρακτηριστικών χειλιών
και παγοποίηση διαχωρισμών train/val/test.
\textbf{Κεφάλαιο~4} εξηγεί τη δημιουργία συνθετικών δεδομένων
(μεταγραφή + TTS για κάθε μηχανή).
\textbf{Κεφάλαιο~5} τεκμηριώνει τις πρώτες δοκιμές
με mel-CNN baseline (PR \#7) και τα αρχικά ευρήματα.
\textbf{Κεφάλαιο~6} περιγράφει τα μοντέλα δεύτερης φάσης:
ηχητικά, οπτικά και late fusion baselines (PRs \#8, \#9).
\textbf{Κεφάλαιο~7} αναλύει το πρόβλημα codec mismatch
(WAV vs.\ MP3) και την αντιμετώπισή του.
\textbf{Κεφάλαιο~8} παρουσιάζει τη στρατηγική voice-disjoint
διαχωρισμού και γιατί είναι αναγκαία.
\textbf{Κεφάλαιο~9} καλύπτει την εκπαίδευση μοντέλων,
την επιλογή υπερπαραμέτρων και τα αποτελέσματα αξιολόγησης.
\textbf{Κεφάλαιο~10} πραγματεύεται τη late fusion:
τι μπορεί και τι δεν μπορεί να κάνει ο concat συνδυασμός,
και γιατί η αυθεντική πολυτροπική συνέπεια παραμένει ανοιχτό ζήτημα.
\textbf{Κεφάλαιο~11} παρουσιάζει την τελική παράδοση:
συγκεντρωτικά αποτελέσματα test split, σύγκριση με παράλληλες
πειραματικές προσεγγίσεις, και ανοιχτά ερωτήματα για μελλοντική εργασία.
